\section{Method}
\label{sec:method}

CrossConfig-CSI converts heterogeneous CSI $\mathbf H$ and acquisition configuration $\mathbf c$
into a structured physical record and an auditable channel report:
\begin{equation}
\hat{\mathbf y}=f_{\Theta}(\mathbf H,\mathbf c),
\qquad
\hat{s}=V\!\left(\mathcal C(\hat{\mathbf y})\right),
\label{eq:method-overview}
\end{equation}
where $\mathcal C$ corrects known physical relations and $V$ deterministically verbalizes the
corrected record.

\begin{figure}[!t]
\centering
\includegraphics[width=0.98\columnwidth]{figures/method_overview_updated.png}
\caption{Overview of \model{}. Stage A phase-aligns, frequency-resamples, and transforms the
heterogeneous CSI into unified beam-frequency tokens while retaining a raw-power bypass. Stage B
uses configuration conditioning and a shared Transformer encoder; text and prototype embeddings
provide semantic supervision only during training. Stage C maps the shared CSI embedding to temporal,
angular, and power/structure predictions, applies relational correction, and deterministically verbalizes
the corrected physical record. Solid arrows denote the test-time inference path, dashed arrows denote
training-only semantic alignment, and the red arrow shows the backpropagation update to attention.}
\label{fig:method-overview}
\end{figure}

Section~III-A introduces the training-only CSI--text--prototype alignment. Section~III-B converts
heterogeneous observations into configuration-conditioned CSI representations. Section~III-C predicts
temporal, angular, power, and multipath attributes with structure-matched branches. Section~III-D
corrects known relations and deterministically verbalizes the resulting record.

\subsection{CSI--Text--Adaptive Prototype Semantic Alignment}
\label{subsec:semantic-alignment}

We organize CSI representations according to physically meaningful propagation states by aligning them
with canonical textual descriptions and adaptive class-level prototypes. Let
$\mathcal K=\{k_m\}_{m=1}^{M}$ denote the semantic classes observed in the training set. Each key
$k_m$ summarizes selected attributes such as environment type, LoS/NLoS condition, multipath richness,
and coarse delay or power regimes. A canonical description $\bar{s}_m$ provides a stable linguistic
definition of each class.

\begin{figure}[!t]
\centering
\includegraphics[width=\columnwidth]{figures/csi_text_prototype_semantic_alignment_replacement.png}
\caption{CSI--text--prototype semantic alignment. The CSI embedding $\mathbf z_c$, frozen textual
anchor $\mathbf z_t$, and adaptive prototype $\mathbf z_p$ are aligned during training so that matched
propagation semantics are close and mismatched classes are separated. At test time, only the CSI encoder
and physics prediction path are used.}
\label{fig:csi-text-prototype-alignment}
\end{figure}

For sample $i$ with semantic label $y_i$, the normalized CSI, text, and prototype representations are
\begin{align}
\mathbf z_c^{\,i}
&=\operatorname{Norm}\!\left(
f_{\mathrm{CSI}}(\mathbf X_i,\mathbf M_i,\mathbf c_i)
\right),
\label{eq:zc-definition}\\
\mathbf z_t^{\,m}
&=\operatorname{Norm}\!\left(f_{\mathrm{text}}(\bar{s}_m)\right),
\qquad
\mathbf z_p^{\,m}=\operatorname{Norm}(\mathbf p_m).
\label{eq:zt-zp-definition}
\end{align}
Here, $\mathbf z_c^{\,i}$ retains sample-specific channel information, whereas $\mathbf z_t^{\,m}$ and
$\mathbf z_p^{\,m}$ describe class $m$. The text encoder is frozen so that the canonical descriptions
remain stable semantic references. The prototype bank is trainable and is initialized from the
corresponding textual anchors,
\begin{equation}
\mathbf p_m^{(0)}\leftarrow\mathbf z_t^{\,m}.
\label{eq:prototype-initialization}
\end{equation}
Thus, each prototype begins with an explicit linguistic meaning but can adapt to the empirical CSI
distribution of its propagation class.

Because all three representations are $L_2$-normalized, their inner products measure cosine similarity.
We use the temperature-scaled similarity
\begin{equation}
s(\mathbf a,\mathbf b)=\exp(\eta)\mathbf a^{\mathsf T}\mathbf b,
\label{eq:semantic-similarity}
\end{equation}
where $\eta$ is a learnable logit-scale parameter. Rather than expanding three nearly identical
cross-entropy expressions, we write the complete alignment objective as
\begin{align}
\mathcal L_{\mathrm{align}}
={}&\lambda_{ct}\mathcal L_{\mathrm{CE}}
\!\left(\mathbf z_c,\{\mathbf z_t^{\,m}\}_{m=1}^{M},y\right)
\nonumber\\
&+\lambda_{cp}\mathcal L_{\mathrm{CE}}
\!\left(\mathbf z_c,\{\mathbf z_p^{\,m}\}_{m=1}^{M},y\right)
\nonumber\\
&+\lambda_{tp}\mathcal L_{\mathrm{NCE}}^{\mathrm{sym}}
\!\left(\{\mathbf z_t^{\,m}\}_{m=1}^{M},
\{\mathbf z_p^{\,m}\}_{m=1}^{M}\right).
\label{eq:semantic-alignment-loss}
\end{align}
The first term moves a CSI sample toward its correct frozen textual anchor, the second adapts the
prototype centers to the CSI distribution, and the third prevents those prototypes from losing their
original physical meanings. We set $\lambda_{ct}=\lambda_{cp}=\lambda_{tp}=1$. The intended relation is
\begin{equation}
\mathbf z_c^{\,i}\approx\mathbf z_t^{\,y_i}\approx\mathbf z_p^{\,y_i},
\label{eq:semantic-relation}
\end{equation}
where $\mathbf z_t$ is fixed, while $\mathbf z_c$ and $\mathbf z_p$ are trainable. This alignment is used
only to shape the CSI embedding space. No text description or prototype lookup is required to generate
a report at inference time.

\subsection{Configuration-Aware Representation Learning}
\label{subsec:configuration-aware}

CSI observations can differ in antenna geometry, frequency resolution, subcarrier spacing, and the
number of valid spatial elements. We therefore construct unified tokens through
\begin{equation}
(\mathbf X,\mathbf M,\mathbf B)=\mathcal T(\mathbf H,\mathbf c),
\label{eq:unified-tokenization}
\end{equation}
where $\mathbf X$ denotes beam--frequency tokens, $\mathbf M$ is their validity mask, and $\mathbf B$
contains normalized beam positions. The operator first removes arbitrary global phase using a reference
coefficient,
\begin{equation}
\mathbf H^{(a)}=\mathbf H\exp\!\left(-j\angle H_{\mathrm{ref}}\right),
\label{eq:phase-alignment}
\end{equation}
resamples the frequency response through the delay domain to a common resolution, and applies an
array-dependent beamspace transform,
\begin{equation}
\mathbf H_{\mathrm{beam}}=
\begin{cases}
\operatorname{FFT}_{\mathrm{ant}}(\mathbf H^{(a)}), & \text{ULA},\\
\operatorname{FFT2}_{\mathrm{row,col}}(\mathbf H^{(a)}), & \text{UPA}.
\end{cases}
\label{eq:beamspace-transform}
\end{equation}
Adjacent ULA beams or local UPA beam patches form tokens by concatenating real and imaginary
components. Invalid array positions and patch boundaries are padded and excluded by $\mathbf M$.

Before entering the shared encoder, only valid token elements are standardized:
\begin{equation}
\widetilde{\mathbf X}=
\frac{\mathbf X-\mu_{\mathbf X}}{\sigma_{\mathbf X}+\epsilon},
\label{eq:masked-standardization}
\end{equation}
where $\mu_{\mathbf X}$ and $\sigma_{\mathbf X}$ are computed over entries selected by $\mathbf M$.
This reduces link-dependent scale variation without allowing padding to affect the statistics.

Each token is processed by a linear projection, a depthwise--pointwise convolution, and adaptive
average/max pooling. The linear layer maps heterogeneous local features to a common hidden dimension,
the convolution captures short-range beam--frequency patterns, and adaptive pooling supports different
input lengths. Subcarrier spacing $\Delta f$ then modulates the projected token through
\begin{equation}
\operatorname{FiLM}(\mathbf u;\Delta f)=
\left(1+\boldsymbol\gamma(\Delta f)\right)\odot\mathbf u
+\boldsymbol\beta(\Delta f).
\label{eq:configuration-film}
\end{equation}
Beam-position and acquisition-configuration encodings are subsequently added:
\begin{equation}
\mathbf e_k=
\operatorname{FiLM}(\mathbf u_k;\Delta f)
+\mathbf e_{\mathrm{beam}}(\mathbf B_k)
+\mathbf e_{\mathrm{cfg}}(\mathbf c).
\label{eq:token-configuration-encoding}
\end{equation}

A shared Transformer processes a learned $[\mathrm{CLS}]$ token together with all valid CSI tokens.
The attention mask prevents padded elements from contributing to contextual aggregation. The final CSI
embedding combines the global $[\mathrm{CLS}]$ state with the masked token mean:
\begin{equation}
\mathbf z_c=\operatorname{Norm}\!\left[
\mathcal P_c\!\left(
\mathbf Z_{\mathrm{cls}}+
\frac{\sum_k M_k\mathbf Z_k}{\sum_k M_k+\epsilon}
\right)\right].
\label{eq:shared-csi-embedding}
\end{equation}
The $[\mathrm{CLS}]$ state captures global interactions learned by attention, whereas the masked mean
preserves evidence distributed across the valid beam--frequency regions. The resulting encoder can share
parameters across configurations without forcing every observation to have the same physical array size
or frequency sampling. During training, the losses in Section~\ref{subsec:semantic-alignment} and
auxiliary semantic and balanced $K$-factor-bin classifiers supervise $\mathbf z_c$; these auxiliary
branches are absent from inference.

\subsection{Physics-Aware Structured Prediction}
\label{subsec:physics-prediction}

The shared embedding is first adapted to physical prediction by a residual MLP:
\begin{equation}
\mathbf z_{\mathrm{phy}}
=\mathbf z_c+\alpha g_{\mathrm{shared}}(\mathbf z_c),
\label{eq:physics-representation}
\end{equation}
where $\alpha$ controls the residual contribution. Separate branches then predict categorical,
temporal, angular, power, and multipath-structure variables, avoiding a single output parameterization
for quantities with incompatible units and geometries.

\paragraph{Temporal prediction.}
Delay-related heads combine $\mathbf z_{\mathrm{phy}}$ with a delay-domain context
\begin{equation}
\mathbf d=g_{\mathrm{delay}}(\mathbf X,\mathbf M,\Delta f).
\label{eq:delay-context}
\end{equation}
They first classify a coarse interval and then regress the normalized position inside that interval:
\begin{align}
\hat b
&=\arg\max_b\operatorname{Softmax}
\!\left(h_{\mathrm{bin}}([\mathbf z_{\mathrm{phy}};\mathbf d])\right)_b,
\label{eq:delay-bin}\\
\hat r
&=\sigma\!\left(h_{\mathrm{pos}}([\mathbf z_{\mathrm{phy}};\mathbf d])\right),
\label{eq:delay-position}\\
\hat y
&=l_{\hat b}+\hat r\left(u_{\hat b}-l_{\hat b}\right),
\label{eq:coarse-to-fine-delay}
\end{align}
where $l_{\hat b}$ and $u_{\hat b}$ are the lower and upper bounds of the selected interval. This
coarse-to-fine construction is more stable than unconstrained regression over a wide delay range.

\paragraph{Angular prediction.}
To respect circular geometry, the angular branch predicts sine and cosine components and reconstructs
the angle as
\begin{equation}
(\hat s_\theta,\hat c_\theta)=h_{\mathrm{ang}}(\mathbf z_{\mathrm{phy}}),
\qquad
\hat\theta=\operatorname{atan2}(\hat s_\theta,\hat c_\theta).
\label{eq:angular-prediction}
\end{equation}
This avoids an artificial discontinuity at the angular wrap-around boundary.

\paragraph{Raw-scale power bypass.}
Standardization removes absolute amplitude, so the unstandardized tokens are processed by a separate
raw-scale encoder:
\begin{equation}
\mathbf r_{\mathrm{raw}}=g_{\mathrm{raw}}(\mathbf X,\mathbf M).
\label{eq:raw-power-feature}
\end{equation}
It summarizes amplitude and power statistics, dominant-beam responses, frequency variation, and
delay-profile energy. In the residual power configuration used here, first-path power is predicted by
\begin{equation}
\hat p_1=p_{\mathrm{base}}(\mathbf z_{\mathrm{phy}})
+g_{\mathrm{res}}([\mathbf z_{\mathrm{phy}};\mathbf r_{\mathrm{raw}}]).
\label{eq:residual-power-prediction}
\end{equation}
The semantic representation determines the propagation regime, while the bypass restores the absolute
scale needed for power estimation.

\paragraph{Multipath structure and training objective.}
Path count, reflection count, delay spread, $K$-factor, and related fields are predicted from suitable
combinations of $\mathbf z_{\mathrm{phy}}$, delay context, raw-scale statistics, and delay-profile
features. Missing annotations are masked so that every sample need not provide every target. Let
$\mathcal Q$ denote the available physical targets. The overall training objective is
\begin{align}
\mathcal L
={}&\sum_{q\in\mathcal Q}\lambda_q\mathcal L_q
+\lambda_{\mathrm{align}}\mathcal L_{\mathrm{align}}
\nonumber\\
&+\lambda_{\mathrm{sem}}\mathcal L_{\mathrm{sem}}
+\lambda_{\mathrm{attr}}\mathcal L_{\mathrm{attr}}
+\lambda_{\mathrm{rel}}\mathcal L_{\mathrm{rel}},
\label{eq:overall-training-objective}
\end{align}
where $\mathcal L_q$ matches the target type, $\mathcal L_{\mathrm{sem}}$ and
$\mathcal L_{\mathrm{attr}}$ provide auxiliary supervision, and $\mathcal L_{\mathrm{rel}}$ softly
penalizes elementary physical inconsistencies. Target-specific weights and delay intervals are provided
in the appendix.

\subsection{Relational Correction and Verbalization}
\label{subsec:correction-verbalization}

Independently predicted fields can violate simple propagation relations. A deterministic correction
operator is therefore applied before text generation:
\begin{equation}
\widetilde{\mathbf y}=\mathcal C(\hat{\mathbf y}).
\label{eq:deterministic-correction}
\end{equation}
It clamps nonnegative quantities, converts count-like outputs into valid integers, and enforces at least
one retained path. Representative corrections are
\begin{align}
\widetilde{\tau}_{\mathrm{RMS}}
&=\max(0,\hat{\tau}_{\mathrm{RMS}}),
\label{eq:spread-correction}\\
\widetilde N_{\mathrm{path}}
&=\max\!\left(1,\operatorname{round}(\hat N_{\mathrm{path}})\right),
\label{eq:path-count-correction}\\
\widetilde N_{\mathrm{ref}}
&=\max\!\left(0,\operatorname{round}(\hat N_{\mathrm{ref}})\right).
\label{eq:reflection-count-correction}
\end{align}
For a predicted LoS sample, the earliest-path delay is reconciled with the dedicated direct-LoS delay:
\begin{equation}
\widetilde\tau_1
=\widetilde\tau_{\mathrm{LoS}}
=\max(0,\hat\tau_{\mathrm{LoS}}).
\label{eq:los-delay-correction}
\end{equation}
Related rules ensure that first-path power does not exceed total power and that mutually incompatible
LoS- and NLoS-specific fields are not emitted together. The operator is non-learned, and both the raw
prediction $\hat{\mathbf y}$ and corrected record $\widetilde{\mathbf y}$ can be retained for auditing.

The deterministic verbalizer produces
\begin{equation}
\hat{s}=V(\widetilde{\mathbf y}).
\label{eq:deterministic-verbalization}
\end{equation}
It selects a conditional template for the predicted propagation state and fills its slots only with
values from the corrected record. Inapplicable fields are suppressed, and qualitative expressions such
as ``short delay,'' ``strong direct component,'' or ``rich multipath'' are determined by predefined
physical intervals. Consequently, every emitted statement can be traced to a predicted field, interval,
or correction rule rather than an unconstrained language-model generation step.

During inference, the complete path is
\begin{equation}
(\mathbf H,\mathbf c)
\rightarrow(\mathbf X,\mathbf M)
\rightarrow\mathbf z_c
\rightarrow\hat{\mathbf y}
\rightarrow\widetilde{\mathbf y}
\rightarrow\hat{s},
\label{eq:complete-test-path}
\end{equation}
with all text, prototype, and auxiliary semantic supervision removed.
